% !TEX program = xelatex
% SPDX-License-Identifier: MIT (templates/latex/LICENSE). 이 템플릿 틀만 MIT 이고, 여기에 쓴 글은 쓴 사람의 것입니다.
% 한국어 논문(짧은 글) 템플릿: XeLaTeX + kotex. 글꼴은 Noto Serif/Sans CJK KR 로 정해 두었다.
\documentclass[11pt,a4paper]{article}

\usepackage[hangul]{kotex}  % hangul: 날짜·장 제목 등을 한국어 방식으로
\setmainhangulfont{Noto Serif CJK KR}[AutoFakeSlant=0.15]  % 한글 기울임꼴이 없어 기울여 만든다
\setsanshangulfont{Noto Sans CJK KR}[AutoFakeSlant=0.15]
\setmonohangulfont{Noto Sans Mono CJK KR}

\usepackage[margin=25mm]{geometry}
\usepackage{amsmath,amssymb}
\usepackage{graphicx}
\usepackage{booktabs}
\usepackage[hidelinks]{hyperref}

% 표·그림 제목 등을 한국어로(article 의 기본값은 "Table", "Figure"). \은 \는 은 kotex 의 조사 자동 선택이다.
\AtBeginDocument{%
  \renewcommand{\tablename}{표}%
  \renewcommand{\figurename}{그림}%
  \renewcommand{\contentsname}{차례}%
  \renewcommand{\refname}{참고 문헌}%
  \renewcommand{\abstractname}{요약}%
}

\title{논문 제목을 여기에 씁니다}
\author{홍길동\\ 소속 기관}
\date{\today}

\begin{document}
\maketitle

\begin{abstract}
이 글은 한국어 LaTeX 문서의 예시다. 요약에는 연구의 목적, 방법, 결과를 짧게 적는다.
\end{abstract}

\section{서론}\label{sec:intro}
한글과 English 를 섞어 쓸 수 있다. 수식은 본문 안에 $E = mc^2$ 처럼 쓰거나 따로 줄을 바꿔 쓴다.
\begin{equation}\label{eq:sum}
  \sum_{k=1}^{n} k = \frac{n(n+1)}{2}
\end{equation}
식~\eqref{eq:sum}\은 자연수의 합이다.

\section{방법}
표~\ref{tab:result}\는 결과를 정리한 것이다.
\begin{table}[h]
  \centering
  \caption{실험 결과}\label{tab:result}
  \begin{tabular}{lrr}
    \toprule
    방법 & 정확도(\%) & 시간(초) \\
    \midrule
    기존 방법 & 81.2 & 12.4 \\
    제안 방법 & 87.5 & 9.8 \\
    \bottomrule
  \end{tabular}
\end{table}

\section{결론}
\ref{sec:intro}절에서 말한 문제를 해결했다. 앞으로 할 일은 다음과 같다.
\begin{itemize}
  \item 더 많은 자료로 실험한다.
  \item 결과를 공개한다.
\end{itemize}

\end{document}
